\documentclass[12pt]{article}
\usepackage[utf8]{inputenc}
\usepackage[T1]{fontenc}
\usepackage{hyperref}
\title{Chemakum Language Resources for K-12 Curriculum Development}
\author{Compiled by Claude Sonnet 5.5 (claude-sonnet-5-5) for Andr\'as Kornai}
\date{}
\begin{document}
\maketitle

\section{Introduction}
Chemakum (ISO 639-3 xch) is an extinct Chimakuan language. It was spoken on the east side of the Olympic Peninsula, between Port Townsend and Hood Canal, until about the 1940s. Franz Boas worked with the last speakers in 1890. According to the sources found, he met only three speakers and collected about 1,250 words, forms and sentences. His main consultant was Louise Webster, whose speech was affected by forgetting parts of the language. Manuel Andrade rechecked some of the material with her in the 1920s. The documentation is very thin. I found no teaching materials, programmes or community revitalization resources.

\section{Primary Educational Resources}
I found no primers, textbooks or classroom materials for Chemakum.

\section{Linguistic and Reference Materials}
\begin{itemize}
\item Wikipedia, ``Chemakum language'': overview and references. \url{https://en.wikipedia.org/wiki/Chemakum_language}
\item Wikipedia, ``Chimakuan languages'': the family, whose other member is Quileute. \url{https://en.wikipedia.org/wiki/Chimakuan_languages}
\item American Philosophical Society, ``Chemakum materials, ACLS collection'': a lexicon of over 1,500 slips compiled by Boas from his 1890 work, a notebook by Manuel Andrade with an English-Chimakum vocabulary, and analyses by Sapir and Swadesh. \url{https://indigenousguide.amphilsoc.org/index.php/entry/40008}
\item A National Library of Medicine digital repository PDF appeared in the results for this search. I did not verify what it contains, so I do not cite it as a Chemakum resource.
\end{itemize}

\section{Religious and Community Resources}
\begin{itemize}
\item Wikipedia, ``Chemakum people'': background on the people. \url{https://en.wikipedia.org/wiki/Chemakum_people}
\end{itemize}
I found no tribal language programme for Chemakum.

\section{Digital and Audio Resources}
I found no audio, video or software resources.

\section{Curriculum Development Context}
There are no speakers and no existing programme. The record is the Boas lexicon and the Andrade vocabulary in the archives. Quileute, the related Chimakuan language, has a revitalization effort, and a project would have to rely on that for comparison. Local history and geography units about the Olympic Peninsula could use the Chemakum people and what remains of their language. Any such unit should say how limited the record is.

\end{document}